\documentclass{article}
\usepackage{geometry}
\geometry{a4paper, margin=1in}
\title{ML Research Proposal: GNN-Based Rule Learning for Bio-Mechanical Cell-to-Cell Communication}
\author{Candidate Name}
\date{Suggested Supervisors: Prof. Dr. Ivo F. Sbalzarini, Dr. Carl D. Modes}

\begin{document}
\maketitle

\section*{Abstract}
The Active Vertex Model in TopoSPAM simulates cell-to-cell communication (e.g., in cell alignment/polarity) using manually defined, simplified neighborhood-averaging rules \cite{1}. Real biological communication, particularly incorporating bioelectric signaling, is a complex, non-linear function of a cell's state and its neighbors' states. This project proposes using \textbf{Graph Neural Networks (GNNs)} to learn the true, complex cell-to-cell communication rule for morphogenesis. We will model the epithelial tissue as a graph, where nodes are cells (with features like $\text{V}_{\text{mem}}$, area $A_a$, and polarity $p_a$) and edges represent cell-to-cell contacts (e.g., gap junctions, adherens junctions). A GNN's message-passing architecture will be trained on high-resolution simulation or experimental data to predict the dynamic update of a cell's polarity and active traction forces based on its neighbors' states. By replacing the simple, hard-coded rules with a learned GNN, we can more accurately model collective cell behavior, infer the underlying bio-mechanical interaction laws, and ultimately generate more predictive TopoSPAM simulations.

\section*{Related Work}
\begin{itemize}
    \item \textbf{Active Vertex Models:} Mechanistic models (e.g., TopoSPAM) that use cell-level topology to simulate tissue dynamics \cite{1}.
    \item \textbf{Graph Neural Networks in Biology:} Emerging use of GNNs for modeling biological networks, including spatial transcriptomics and protein-protein interactions (e.g., ScaDS.AI/Computer Science expertise).
    \item \textbf{Bioelectric Networks:} The use of bioelectrical signaling ($\text{V}_{\text{mem}}$) via gap junctions as the communication channel to be learned by the network (Levin, 2012; Barriga Lab data).
\end{itemize}

\end{document}
